% translation-source: docs/user/latex/chapters/05-lumiere.tex; sha256: 20046c985a406213c2f06ec61857c03dcac1c7cd28fcd7f27adae0d8386db34a
\chapter{Light and histogram}
\section{Read the fixed histogram}
The histogram stays at the top of Adjustments while scrolling. It describes
the processed current sRGB preview, not raw sensor values. End indicators
flag preview values at their limits; they do not directly measure sensor
saturation.

Reading zones are Blacks(0--10\,\%), Shadows(10--33\,\%),
Exposure(33--66\,\%), Highlights(66--90\,\%) and Whites(90--100\,\%).
These are visual references, not hard slider influence boundaries.

\section{Controls}
\begin{longtable}{@{}L{3cm}L{2.8cm}L{7.8cm}@{}}
\toprule Setting & Range & Usage\\\midrule\endhead
Exposure & $-5$ to $+5$ EV & Adjust overall brightness; begin with small changes.\\
Contrast & $-100$ to $+100$ & Strengthen or soften tonal separation.\\
Shadows & $-100$ to $+100$ & Brighten or darken dark regions, not just the black endpoint.\\
Blacks & $-100$ to $+100$ & Adjust the darkest tones and the image's tonal foundation.\\
Highlights & $-100$ to $+100$ & Moderate or strengthen bright tones.\\
Whites & $-100$ to $+100$ & Adjust tones nearest white.\\
\bottomrule
\end{longtable}
Zero is the reference for all six controls. On a RAW with camera signal and
DCP, exposure operates before the profile: $+1$EV doubles the linear signal
there, not necessarily TIFF values after the curve. Without that camera
context, the gain is applied to the linear master. The \emph{Grayscale}
checkbox in this group toggles monochrome rendering.

\section{Suggested procedure}
Set exposure first, then shadows and highlights, followed by blacks, whites
and contrast. Revisit earlier steps as needed: these controls interact
visually. Raising shadows also reveals existing noise. A highlight adjustment
cannot recreate information absent from the source.

\paragraph{Example.}
For backlighting, try slightly lowering highlights and moderately raising
shadows rather than immediately pushing both sliders to their extremes.
Hold \touche{Space} in the image view to compare, then inspect transitions
using \touche{Ctrl + 1}.
\gestesreglages
